%%%%%%%%%%%%%%%%%%%%%%%%%%%%%%%%%%%%%%%
% Cover Letter - BCG, Working Student Generative AI & Workflow Automation
%
% Compile with: cd cover_letters && xelatex cover_bcg_genai.tex
%%%%%%%%%%%%%%%%%%%%%%%%%%%%%%%%%%%%%%%

\documentclass[]{cover}
\usepackage{fancyhdr}
\pagestyle{fancy}
\fancyhf{}
\rfoot{Page \thepage \hspace{0pt}}
\thispagestyle{empty}
\renewcommand{\headrulewidth}{0pt}
\begin{document}

%%%%%%%%%%%%%%%%%%%%%%%%%%%%%%%%%%%%%%
%     TITLE NAME
%%%%%%%%%%%%%%%%%%%%%%%%%%%%%%%%%%%%%%

\namesection{}{\Huge{Dalfin Rafael Solihin}}{
  \href{mailto:dalfinrafa@gmail.com}{dalfinrafa@gmail.com} | +49 155 664 752 60 | \urlstyle{same}\href{https://www.linkedin.com/in/dalfin-rafael-solihin}{LinkedIn}
}

%%%%%%%%%%%%%%%%%%%%%%%%%%%%%%%%%%%%%%
%     MAIN COVER LETTER CONTENT
%%%%%%%%%%%%%%%%%%%%%%%%%%%%%%%%%%%%%%

\currentdate{\today}

\lettercontent{Sehr geehrtes BCG-Team,}

\lettercontent{ich bewerbe mich für die Stelle als Working Student Generative AI \& Workflow Automation am Münchner Standort. Ich habe bereits ein vollständiges LLM-Automatisierungs-Workflow gebaut und möchte genau diese Erfahrung nun einsetzen, um BCG-Teams dabei zu unterstützen, ähnliche KI-Lösungen zu entwickeln und produktiv zu nutzen.}

\lettercontent{Mein direktester Beitrag für diese Rolle ist ein konkretes Projekt: Ich habe einen AI Application Generator in Python entwickelt, der über die Google Gemini API Bewerbungsunterlagen vollständig automatisiert generiert. Das System umfasst API-Integration, iteratives Prompt-Engineering für konsistente Ausgabequalität sowie eine sichere Datenpipeline mit PII-Stripping und automatisierter PDF-Erzeugung --- entwickelt mit Claude Code als agentischem Entwicklungswerkzeug, das ich täglich einsetze. Ergänzend dazu habe ich im Rahmen meines Forschungsassistentenjobs an der HM München Python-Skripte zur Auswertung von SFRA-Messdaten entwickelt und Ergebnisse für ein Publikum mit unterschiedlichem Fachkenntnisstand dokumentiert. Drei Fähigkeiten, die direkt auf die ausgeschriebenen Aufgaben einzahlen:}

{\raggedright\fontspec[Path = OpenFonts/fonts/raleway/]{Raleway-Medium}\fontsize{11pt}{13pt}\selectfont
\begin{itemize}
    \item \textbf{API-Integration und Prompt-Engineering:} End-to-End-Implementierung mit Google Gemini API, sicherer Datenpipeline und iterativem Prompt-Design für produktionsreife Ausgabequalität
    \item \textbf{KI-Tool-Evaluation:} Training und quantitative Auswertung eines YOLOv8-Bildklassifikationsmodells sowie laufende Erprobung neuer LLM-Features im AI Application Generator
    \item \textbf{Technische Kommunikation:} Übersetzung komplexer Messergebnisse in strukturierte Berichte für Stakeholder mit und ohne technischen Hintergrund --- Grundlage für Schulungen und Enablement-Formate
\end{itemize}\par}
\vspace{6pt}

\lettercontent{BCGs 10-20-70-Prinzip --- 10\,\% Algorithmen, 20\,\% Daten und Technologie, 70\,\% Menschen und Prozesse --- beschreibt genau den Schwerpunkt dieser Stelle: nicht nur Tools bauen, sondern Teams dazu befähigen, sie zu nutzen. BCG X zeigt, dass BCG an der Grenze des Möglichen in der KI-Produktentwicklung arbeitet, und das Center for Responsible Generative AI unterstreicht, dass Verantwortungsbewusstsein dabei kein Nachgedanke ist. Dass Anthropic --- der Hersteller von Claude Code, meinem täglichen Entwicklungswerkzeug --- ein namentlich genannter BCG-Partnersystempartner ist, schließt den Kreis.}

\lettercontent{Ich lerne neue Tools schnell, arbeite sehr selbstständig und bin in der Lage, technische Konzepte klar zu erklären --- wie mein Arbeitszeugnis von der HM München belegt. Die Kombination aus praktischer KI-Erfahrung und der Fähigkeit, Brücken zwischen Technik und Nicht-Technik zu bauen, ist genau das, was eine Enablement-Rolle wie diese verlangt.}

\lettercontent{Ich freue mich darauf, von Ihnen zu hören.}

\begin{flushright}
\closing{Mit freundlichen Grüßen,}
\signature{Dalfin Rafael Solihin}
\end{flushright}

\end{document}
